% ===== PRÉAMBULE COMMUN — À COPIER TEL QUEL EN TÊTE DE CHAQUE FICHIER .tex =====
% Ne modifier QUE :
%   - les deux lignes \fancyhead[L] et \fancyhead[R]
%   - le bloc titre \begin{center}...\end{center} juste après \begin{document}
\documentclass[11pt,a4paper]{article}
\usepackage[utf8]{inputenc}
\usepackage[T1]{fontenc}
\usepackage{lmodern}
\usepackage[french]{babel}
\usepackage{amsmath,amssymb,mathtools}
\usepackage{icomma}
\usepackage{xcolor}
\usepackage{tikz}
\usepackage{pgfplots}
\usepackage[most]{tcolorbox}
\usepackage{enumitem}
\usepackage{array,booktabs,colortbl}
\usepackage[a4paper,margin=2.2cm,headheight=26pt,headsep=10pt]{geometry}
\usepackage{fancyhdr}
\usepackage[hidelinks]{hyperref}
\usetikzlibrary{arrows.meta,calc,angles,quotes,positioning,decorations.markings,intersections,backgrounds}
\pgfplotsset{compat=1.18}

\definecolor{primary}{RGB}{226,74,88}
\definecolor{primarydark}{RGB}{150,28,42}
\definecolor{methcol}{RGB}{40,90,150}
\definecolor{excol}{RGB}{0,135,90}
\definecolor{defcol}{RGB}{0,114,178}
\definecolor{attncol}{RGB}{200,40,55}
\definecolor{thmcol}{RGB}{0,140,120}
\definecolor{courbe}{RGB}{32,110,200}
\definecolor{courbeder}{RGB}{210,70,40}
\definecolor{tang}{RGB}{0,150,90}
\definecolor{grille}{RGB}{200,205,215}
\definecolor{plus}{RGB}{200,60,60}
\definecolor{moins}{RGB}{30,90,200}

\tcbset{boxbase/.style={enhanced,breakable,boxrule=0.9pt,arc=2.5pt,
  left=3mm,right=3mm,top=2.3mm,bottom=2.3mm,fonttitle=\bfseries,coltitle=white}}
\newtcolorbox{cle}[1][]{boxbase,colback=defcol!6,colframe=defcol,
  title={\textsc{À retenir}\ifx\relax#1\relax\relax\else\textmd{ : \itshape #1}\fi}}
\newtcolorbox{methode}[1][]{boxbase,colback=methcol!7,colframe=methcol,sharp corners,
  title={\textsc{Méthode}\ifx\relax#1\relax\relax\else\textmd{ --- \itshape #1}\fi}}
\newtcolorbox{exemplebox}[1][]{boxbase,colback=excol!5,colframe=excol,
  title={\textsc{Exemple}\ifx\relax#1\relax\relax\else\textmd{ : \itshape #1}\fi}}
\newtcolorbox{piege}[1][]{boxbase,colback=attncol!6,colframe=attncol,
  title={\textsc{Attention}\ifx\relax#1\relax\relax\else\textmd{ : \itshape #1}\fi}}
\newtcolorbox{exo}[1][]{boxbase,colback=primary!4,colframe=primary,
  title={\textsc{Exercice}\ifx\relax#1\relax\relax\else\textmd{~#1}\fi}}
\newtcolorbox{corr}[1][]{boxbase,colback=thmcol!5,colframe=thmcol,
  title={\textsc{Corrigé}\ifx\relax#1\relax\relax\else\textmd{~#1}\fi}}

\newcommand{\R}{\mathbb{R}}
\newcommand{\N}{\mathbb{N}}
\newcommand{\ds}{\displaystyle}
\newcommand{\tg}{\operatorname{tg}}
\newcommand{\cotg}{\operatorname{cotg}}
\newcommand{\arctg}{\operatorname{arctg}}
\newcommand{\arccotg}{\operatorname{arccotg}}
\newcommand{\limh}{\lim_{h\to 0}}

\pagestyle{fancy}\fancyhf{}
\fancyhead[L]{\small\textcolor{primarydark}{\textbf{Thème 2} — Formulaire de dérivation}}
\fancyhead[R]{\small\textcolor{primarydark}{\textsc{Tutoriel}}}
\fancyfoot[C]{\small\thepage}
\renewcommand{\headrulewidth}{0.8pt}

\begin{document}

\begin{center}
{\LARGE\bfseries\color{primarydark} Formulaire de dérivation}\\[2pt]
{\color{primary}\rule{0.5\textwidth}{1.2pt}}\\[4pt]
{\small\itshape Règles directes : constante, puissance, somme, produit et quotient}
\end{center}
\vspace{2mm}

\begin{cle}[les six règles de base]
Pour deux fonctions dérivables $U$ et $V$ et une constante $a\in\R$ :
\begin{center}
\renewcommand{\arraystretch}{1.5}
\begin{tabular}{@{}c l >{$\ds}c<{$}@{}}
\toprule
\textbf{N\textsuperscript{o}} & \textbf{Fonction} & \textbf{Dérivée} \\
\midrule
1 & Constante $\;y=C$ & C'=0 \\
2 & Variable $\;y=x$ & x'=1 \\
3 & Puissance $\;y=x^{n}$ & (x^{n})'=n\,x^{\,n-1} \\
4 & Somme et multiple $\;y=aU+bV$ & (aU+bV)'=aU'+bV' \\
5 & Produit $\;y=UV$ & (UV)'=U'V+UV' \\
6 & Quotient $\;y=\dfrac{U}{V}$ & \left(\frac{U}{V}\right)'=\frac{U'V-UV'}{V^{2}} \\
\bottomrule
\end{tabular}
\end{center}
La règle 3 est valable pour \emph{tout} exposant réel $n$ (entier, négatif ou fractionnaire) : c'est elle qui gère les racines et les inverses.
\end{cle}

\begin{methode}[réécrire les racines et les inverses en puissances]
Avant de dériver, on transforme toute racine ou tout inverse en une puissance $x^{n}$, puis on applique $(x^{n})'=n\,x^{n-1}$ :
\[
\sqrt{x}=x^{1/2},\qquad \sqrt[3]{x^{2}}=x^{2/3},\qquad \frac{1}{x^{n}}=x^{-n},\qquad \frac{1}{\sqrt{x}}=x^{-1/2}.
\]
\textbf{Exemple d'automatisme :} $y=\sqrt{x}=x^{1/2}\ \Rightarrow\ y'=\tfrac12\,x^{-1/2}=\dfrac{1}{2\sqrt{x}}$.\\
De même $y=\dfrac{1}{x^{3}}=x^{-3}\ \Rightarrow\ y'=-3\,x^{-4}=-\dfrac{3}{x^{4}}$.
\end{methode}

\begin{piege}[les deux erreurs classiques]
\begin{itemize}[leftmargin=1.4em,itemsep=1pt]
\item \textbf{Produit :} $(UV)'\neq U'V'$. Il faut \emph{deux} termes : $(UV)'=U'V+UV'$.
\item \textbf{Quotient :} l'ordre du numérateur est \emph{imposé} : $\left(\dfrac{U}{V}\right)'=\dfrac{U'V-UV'}{V^{2}}$, \textbf{jamais} $\dfrac{UV'-U'V}{V^{2}}$. Le terme qui commence par $U'$ vient en premier, et le dénominateur est $V^{2}$ (le carré du dénominateur, pas de $U$).
\end{itemize}
\end{piege}

\begin{exemplebox}[un produit : $y=x\sqrt{x}$]
On pose $U=x$ et $V=\sqrt{x}=x^{1/2}$, d'où $U'=1$ et $V'=\tfrac12 x^{-1/2}=\dfrac{1}{2\sqrt{x}}$.
\[
y'=U'V+UV'=1\cdot\sqrt{x}+x\cdot\frac{1}{2\sqrt{x}}
=\sqrt{x}+\frac{x}{2\sqrt{x}}=\sqrt{x}+\frac{\sqrt{x}}{2}=\frac{3}{2}\sqrt{x}.
\]
\textbf{Vérification} par réécriture directe : $y=x\cdot x^{1/2}=x^{3/2}$, donc $y'=\tfrac32 x^{1/2}=\tfrac32\sqrt{x}$. \checkmark
\end{exemplebox}

\begin{exemplebox}[un quotient : $y=\dfrac{x^{2}}{x^{2}+1}$]
On pose $U=x^{2}$ et $V=x^{2}+1$, d'où $U'=2x$ et $V'=2x$.
\[
y'=\frac{U'V-UV'}{V^{2}}
=\frac{2x\,(x^{2}+1)-x^{2}\cdot 2x}{(x^{2}+1)^{2}}
=\frac{2x^{3}+2x-2x^{3}}{(x^{2}+1)^{2}}
=\frac{2x}{(x^{2}+1)^{2}}.
\]
Le terme $U'V$ est bien écrit \emph{en premier}, et on développe le numérateur avant de simplifier.
\end{exemplebox}

\end{document}
